
%%%%%%%%%%%%%%%%%%%%%%%%%%%%%%%%%%%%%%%%%%%%%%%%%%%%%%%%%%%%%%%%%%%%%%%%%%%%%%%%
% Template for USENIX papers.
%
% History:
%
% - TEMPLATE for Usenix papers, specifically to meet requirements of
%   USENIX '05. originally a template for producing IEEE-format
%   articles using LaTeX. written by Matthew Ward, CS Department,
%   Worcester Polytechnic Institute. adapted by David Beazley for his
%   excellent SWIG paper in Proceedings, Tcl 96. turned into a
%   smartass generic template by De Clarke, with thanks to both the
%   above pioneers. Use at your own risk. Complaints to /dev/null.
%   Make it two column with no page numbering, default is 10 point.
%
% - Munged by Fred Douglis <douglis@research.att.com> 10/97 to
%   separate the .sty file from the LaTeX source template, so that
%   people can more easily include the .sty file into an existing
%   document. Also changed to more closely follow the style guidelines
%   as represented by the Word sample file.
%
% - Note that since 2010, USENIX does not require endnotes. If you
%   want foot of page notes, don't include the endnotes package in the
%   usepackage command, below.
% - This version uses the latex2e styles, not the very ancient 2.09
%   stuff.
%
% - Updated July 2018: Text block size changed from 6.5" to 7"
%
% - Updated Dec 2018 for ATC'19:
%
%   * Revised text to pass HotCRP's auto-formatting check, with
%     hotcrp.settings.submission_form.body_font_size=10pt, and
%     hotcrp.settings.submission_form.line_height=12pt
%
%   * Switched from \endnote-s to \footnote-s to match Usenix's policy.
%
%   * \section* => \begin{abstract} ... \end{abstract}
%
%   * Make template self-contained in terms of bibtex entires, to allow
%     this file to be compiled. (And changing refs style to 'plain'.)
%
%   * Make template self-contained in terms of figures, to
%     allow this file to be compiled. 
%
%   * Added packages for hyperref, embedding fonts, and improving
%     appearance.
%   
%   * Removed outdated text.
%
%%%%%%%%%%%%%%%%%%%%%%%%%%%%%%%%%%%%%%%%%%%%%%%%%%%%%%%%%%%%%%%%%%%%%%%%%%%%%%%%

\documentclass[letterpaper,twocolumn,10pt]{article}
\usepackage{specific/usenix}

% Math packages
\usepackage{amsmath,amssymb}

\usepackage[most]{tcolorbox}
\usepackage{tikz}
\newcommand*\circled[1]{\tikz[baseline=(char.base)]{
             \node[shape=circle,draw=red,inner sep=.5pt] (char) {\textcolor{red}{#1}};}}
             
\lstdefinestyle{mystyle}{
%  backgroundcolor=\color{lightgray},
  commentstyle=\color{codegreen},
  keywordstyle=\color{magenta},
%  numberstyle=\small\color{codegray},
  stringstyle=\color{codepurple},
  basicstyle=\ttfamily\scriptsize,
  breakatwhitespace=false,
  breaklines=true,
  captionpos=b,
  keepspaces=true,
%  numbers=left,
%  numbersep=3pt,
  showspaces=false,
  showstringspaces=false,
  showtabs=false,
  tabsize=2
}


% ---------- Common, venue-agnostic preamble ----------
\usepackage[utf8]{inputenc}
\usepackage[english]{babel}

% Math / graphics
%\usepackage{amsmath,amssymb}
\usepackage{graphicx}
\DeclareGraphicsExtensions{.pdf,.png,.jpg}

% Tables / lists
\usepackage{booktabs}
\usepackage{multirow}
\usepackage{enumitem}
\usepackage{makecell}
\usepackage{relsize}

% Colors & drawings & boxes
\usepackage{xcolor}
\usepackage[most]{tcolorbox}
\usepackage[normalem]{ulem}
\usepackage{colortbl}

% Listings (portable)
\definecolor{darkgreen}{RGB}{0,100,0}
\definecolor{codegreen}{rgb}{0,0.5,0}
\definecolor{codepurple}{rgb}{0.58,0,0.82}
\definecolor{codegray}{rgb}{0.5,0.5,0.5}
\usepackage{listings}
\lstdefinestyle{mystyle}{
  commentstyle=\color{codegreen},
  keywordstyle=\bfseries,
  stringstyle=\color{codepurple},
  basicstyle=\ttfamily\scriptsize,
  breaklines=true,
  captionpos=b,
  keepspaces=true,
  tabsize=2
}
\lstset{style=mystyle}

% Cross-refs (safe as long as we don't load 'caption' ourselves)
%\usepackage[capitalize,nameinlink]{cleveref}

% Algorithms (you can drop this if unused)
\usepackage[noend,ruled,linesnumbered]{algorithm2e}
%\usepackage[ruled,vlined,linesnumbered]{algorithm2e}

% Hyphenation
\hyphenation{op-tical net-works semi-conduc-tor}

\usepackage{tikz}
\usetikzlibrary{arrows.meta,positioning,shapes,calc} % needed by the figure

% ---------- Common macros ----------
\newcommand{\find}[1]{%
\begin{tcolorbox}[tile,size=fbox,boxsep=2mm,boxrule=0pt,top=0pt,bottom=0pt,
borderline={0.6mm}{0pt}{black!66!white},colback=black!5!white]
\em #1
\end{tcolorbox}
}

\newcommand{\question}[1]{%
\begin{tcolorbox}[tile,size=fbox,boxsep=1.5mm,boxrule=0pt,top=0pt,bottom=0pt,
borderline west={1mm}{0pt}{blue!5!white},colback=blue!5!white]
\em #1
\end{tcolorbox}
}

\newcommand{\modify}[1]{\textcolor{blue}{#1}}
\newcommand{\changemark}[1]{{\color{black}#1}}
\newcommand{\redout}[1]{\sout{#1}}

\newcommand{\Honggfuzz}{\href{https://github.com/google/honggfuzz}{Honggfuzz}}
\newcommand{\AFLpp}{\href{https://github.com/AFLplusplus/AFLplusplus}{AFL++}}

\newcommand{\firstAuthor}{$^{\pdftooltip{\textrm{\faUser}}{first author}}$}
\newcommand{\corrAuthor}{$^{\pdftooltip{\textrm{\faEnvelope}}{corresponding author}}$}

\newcommand{\denseitems}{\setlength{\itemsep}{1pt}\setlength{\parskip}{0pt}}

\newcommand{\tclcolor}{\cellcolor[HTML]{C0C0C0}}

\newboolean{COMMENTSON} 
\setboolean{COMMENTSON}{true}   
\ifthenelse{\boolean{COMMENTSON}}
{
\newcommand{\WL}[2]{\todo[color=red!40,linecolor=red,bordercolor=red]{Wen: #2 } #1}
\newcommand{\wen}[1]{\todo[inline,color=red!40,linecolor=red,bordercolor=red]{Wen: #1 }}
}

\definecolor{DarkOrange}{rgb}{0.8,0.3,0.0} 
\definecolor{DarkCyan}{rgb}{0.0, 0.55, 0.55}
\definecolor{codegreen}{rgb}{0,0.6,0}
\definecolor{codegray}{rgb}{0.5,0.5,0.5}
\definecolor{codepurple}{rgb}{0.58,0,0.82}
\definecolor{backcolour}{rgb}{0.95,0.95,0.92}


% ---- Paper identity ----
\newcommand{\papertitle}{Beyond Source: An Empirical Study of Python Bytecode Security Risks}
\newcommand{\papertitleshort}{Python Bytecode Security Risks}

\newcommand{\tech}{\mbox{\textsc{PyBC-Sec}}}
\newcommand{\tool}{\tech}

% ---- Keywords ----
\newcommand{\paperkeywords}{Python bytecode, package security, software supply chain, program analysis, runtime robustness}

% ---- Authors ----
\newcommand{\authorAname}{Wen Li}
\newcommand{\authorAaffil}{Utah State University}
\newcommand{\authorAemail}{awen.li@usu.edu}

\newcommand{\authorBname}{Coauthor Name}
\newcommand{\authorBaffil}{Somewhere University}
\newcommand{\authorBemail}{coauthor@example.com}

% appendix for journal
\newcommand{\appendixcontent}{%
  \appendix
}


%-------------------------------------------------------------------------------
\begin{document}
%-------------------------------------------------------------------------------

%don't want date printed
\date{}

% make title bold and 14 pt font (Latex default is non-bold, 16 pt)
\title{\Large \bf {\tech}: Your Prototype's Name}

%for single author (just remove % characters)
%\author{
%{\rm Baihong Chen~~\firstAuthor}\\
%Utah State University
%\and
%{\rm Yongxi Sha}\\
%Utah State University
%\and
%{\rm Wei Wei}\\
%Utah State University
%\and
%{\rm Wen Li~~\corrAuthor}\\
%Utah State University
%} % end author

\maketitle

%!TEX root = paper.tex

\begin{abstract}
Python security analysis is largely source-centric: package scanners, dependency analyzers, malware detectors, and program-analysis pipelines typically inspect Python source files, metadata, installation scripts, and dependency relations.  However, CPython executes compiled code objects, routinely materializes bytecode as \texttt{.pyc} files, and can load bytecode through source-less modules, cached imports, bundled applications, custom loaders, and marshalled code objects.  This creates a practical visibility gap: behavior may be present in executable bytecode even when corresponding source is missing, stale, mismatched, or difficult to recover.

This paper presents an empirical study of Python bytecode as a first-class security artifact.  We study how bytecode appears in real software distributions, whether it is visible from source, how well existing tools recover source-level representations, and how CPython behaves when processing unusual bytecode inputs.  Our methodology combines package-scale artifact extraction, source/bytecode correspondence analysis, decompilation and recompilation checks, runtime robustness testing across CPython versions, and source-reachability analysis for abnormal behaviors.  The study is designed to separate ecosystem visibility risks from interpreter robustness findings and to distinguish bytecode-only behavior from behavior that ordinary source-level tooling can observe.

Our results show that Python bytecode deserves explicit treatment in package security workflows.  Bytecode artifacts occur in multiple distribution forms, source correspondence is not guaranteed, decompilation-based recovery is incomplete, and bytecode-level testing reaches runtime behaviors that are not always reproducible from ordinary source.  These findings motivate bytecode-aware package scanning, explicit source/bytecode mismatch checks, more robust bytecode validation boundaries, and security tooling that treats compiled Python artifacts as analyzable inputs rather than disposable caches.

\end{abstract}

\section{Introduction}\label{sec:intro}

\begin{itemize}[leftmargin=1.25em]
  \item \textbf{Paragraph 1:} Introduce the broader context and importance.
  \item \textbf{Paragraph 2:} Describe specific limitations of current approaches.
  \item \textbf{Paragraph 3:} Present the key idea and brief intuition.
  \item \textbf{Paragraph 4:} List main contributions in bullet form.
  \item \textbf{Paragraph 5:} Summarize empirical results and impact.
\end{itemize}

\vspace{3pt}
\noindent
\textbf{Open Science.}
The source code of {\tech}, together with all experimental datasets, is publicly available in our 
\href{https://figshare.com/s/344b026e85fe121afbd5}{\underline{artifact package}}. 
All artifacts have been released to support full transparency and reproducibility~\cite{ossfuzz}.

\section{Background and Motivation}\label{sec:background}



\subsection{Key Concepts and Techniques}



\subsection{Motivating Example}\label{sec:motiexample}


\section{The {\tech} Framework}\label{sec:tech}
\begin{itemize}[leftmargin=1.25em]
  \item \textbf{Design Goals:} List desired properties or constraints.
  \item \textbf{High-Level Architecture:} Provide a diagram and describe components.
  \item \textbf{Core Techniques:}
    \begin{itemize}
      \item State the goal of each module.
      \item Explain the key idea or intuition.
      \item Provide detailed step-by-step methodology.
      \item Optionally include a small illustrative example.
    \end{itemize}
  \item \textbf{Implementation Choices:} Note major engineering decisions and tradeoffs.
\end{itemize}

\subsection{Framework Overview}\label{ss:overview}

%!TEX root = paper.tex
\subsection{Dataset Construction and Artifact Extraction}\label{sec:method:dataset}

We construct datasets from multiple sources to reduce bias.  The first dataset contains popular PyPI packages selected by download counts or dependency centrality.  The second contains randomly sampled PyPI packages across popularity and age.  The third contains packages reported as malicious or suspicious by public advisories, security vendors, or package-index takedowns.  The fourth contains Python application bundles and release artifacts such as wheels, archives, container layers, and PyInstaller-style outputs.  The fifth contains runtime-testing inputs derived from CPython tests, real bytecode artifacts, and generated bytecode variants.

For each package or artifact, {\tech} records the package name and version, distribution type, file inventory, source and bytecode paths, bytecode magic number, inferred Python version, hashes, and import or load behavior when such behavior can be determined statically.  The extractor parses \texttt{.pyc} headers, recovers marshalled code objects, disassembles instruction streams, recursively visits nested code objects, and records metadata such as constants, names, variables, flags, exception tables, closures, and line-position information.  It also scans source for dynamic loading APIs including \texttt{marshal}, \texttt{importlib}, custom loaders, and direct calls to \texttt{exec} over recovered code objects.

This design addresses the input-soundness question in two ways.  Multiple dataset classes reduce dependence on one distribution channel, while artifact hashes and explicit version metadata make every extracted input reproducible.  Ambiguous cases are not discarded silently: unknown bytecode versions, unparsable artifacts, encrypted or packed payloads, and dynamic loaders whose target cannot be resolved are reported as separate categories in the evaluation.

\subsection{Phase 2: Practical Analyzability}\label{sec:method:source}

The second phase measures the practical boundary of bytecode analysis.  The unit is a \texttt{.pyc} file, not a package: for each bytecode-containing artifact, the analyzer extracts every \texttt{.pyc} entry, resolves its CPython tag, and tests whether selected tools can load, disassemble, or decompile it.  The outcome is the highest level reached: not loadable, loadable only, disassemblable, partially decompilable, or fully decompilable.

The phase is version-aware by design.  Python bytecode is an implementation detail that can change across CPython releases, and the standard disassembler documents this version sensitivity~\cite{python-dis}.  Therefore, a file is included in tool analysis only when the scan identifies a supported modern CPython tag and a matching interpreter is prepared.  In the current scope, this means CPython~3.8 and later; CPython~3.7 and earlier, PyPy tags, unknown tags, and noncanonical tags without an interpreter mapping remain in the raw scan inventory but are excluded from tool-failure counts.  This prevents missing interpreters or unsupported runtimes from being misreported as tool limitations.

The environment-preparation step reads the scanner's version summary, uses the interpreter field to identify canonical CPython runtimes, locates or installs matching interpreters when the host supports it, and prepares per-version virtual environments for tools that depend on bytecode format.  Each analysis step records both status and failure reason, including marshal failure, disassembly failure, unsupported bytecode version, external-tool timeout, nonzero exit, and unavailable decompiler output.  These failure reasons are necessary for interpretation: an environment miss, a loader failure, and a decompiler failure imply different limits.

Tool selection follows five criteria.  \textbf{C1: Availability} requires public, locally runnable tools so the study does not depend on closed services.  \textbf{C2: Python~3 support} keeps the baseline aligned with the modern CPython bytecode observed in the scan.  \textbf{C3: Practical use} favors tools that analysts or developers are likely to encounter through documentation, packaging, or community references.  \textbf{C4: Maintenance or stability} avoids failures caused mainly by obsolete or unreproducible tooling.  \textbf{C5: Method diversity} ensures that the selected set covers loading, disassembly, and decompilation rather than one tool family.  These criteria define a reproducible baseline for this study; they do not rank all Python bytecode tools globally.

\begin{table}[t]
  \centering
  \caption{Selected bytecode-analysis tools and selection criteria.}
  \label{tab:tool-selection}
  \begin{tabular}{@{}llp{0.25\linewidth}p{0.35\linewidth}@{}}
    \toprule
    Tool & Criteria & Functionality & Role \\
    \midrule
    CPython \texttt{marshal}~\cite{python-marshal} & C1,C2,C4,C5 & Loads code objects from \texttt{.pyc} payloads. & Reference code-object loading. \\
    CPython \texttt{dis}~\cite{python-dis} & C1,C2,C4,C5 & Disassembles bytecode instructions. & Standard bytecode disassembly. \\
    \texttt{uncompyle6}~\cite{uncompyle6} & C1--C5 & Decompiles Python bytecode to source. & Established decompilation baseline. \\
    \texttt{decompyle3}~\cite{decompyle3} & C1--C5 & Python~3 decompilation. & Python-3-focused decompiler baseline. \\
    \texttt{pycdc}~\cite{pycdc} & C1--C5 & Disassembles and decompiles bytecode in C++. & Independent C++ decompiler/disassembler. \\
    PyLingual~\cite{pylingual} & C1--C5 & Modern CPython decompilation. & Recent Python~3 decompiler. \\
    \bottomrule
  \end{tabular}
\end{table}

Table~\ref{tab:tool-selection} applies these criteria to the measurement stack.  CPython \texttt{marshal} and \texttt{dis} provide reference loading and disassembly baselines; the decompilers test whether bytecode can be lifted back toward source-level workflows.  We do not claim this set exhausts all Python bytecode tooling.  It is selected to cover the main local analysis levels available to users while avoiding overclaiming about tools not evaluated.

\subsection{Runtime Robustness and Source Reachability}\label{sec:method:runtime}

To study runtime robustness, {\tech} executes bytecode-level inputs under multiple CPython versions, including CPython 3.10, 3.11, and 3.12, and optionally debug or sanitizer builds.  Inputs include real bytecode artifacts, structured bytecode variants, and a raw byte-mutation baseline.  Outcomes are classified as load failure, Python-level exception, normal termination, timeout, assertion failure, abort, segmentation fault, or sanitizer finding.  We deduplicate abnormal outcomes using signal type, top stack frame, failure location, assertion message, crash-address bucket, input hash, and code-object feature signatures.

For each abnormal outcome, {\tech} checks source reachability.  It attempts to decompile the bytecode, recompile the recovered source, re-execute under the same runtime, and compare the failure signature.  Outcomes are classified as source-reproducible, bytecode-only under current tools, decompilation failed, recompilation failed, behavior mismatch, or inconclusive.  This step separates bytecode-level robustness findings from ordinary source-level vulnerabilities and identifies where source-based testing cannot reproduce the behavior observed at the bytecode boundary.

This stage is included because robustness results without reachability are easy to overstate.  Raw bytecode mutation can generate artifacts that no normal compiler would emit, while ordinary source-level testing may miss behaviors that valid loaders can still execute.  Reporting source reachability alongside crash class keeps these interpretations separate.



\section{Implementation and Limitations} \label{sec:impl}


\subsection{Key Component Implementation} \label{sec:impl:components}


\subsection{Limitations} \label{sec:impl:limits}


\section{Results}\label{sec:results}

This section reports the study results by research question.  RQ1 uses the completed PyPI scan; later result tables are kept as placeholders until the remaining measurements are inserted.  The reporting structure matches the four goals in Section~\ref{sec:study} and the methodology in Section~\ref{sec:method}.

\noindent
{\bf Experiment Environment.}
Experiments are conducted in isolated Linux environments.  Package collection downloads public PyPI artifacts; later analysis phases run locally with explicit timeouts and version-specific CPython interpreters.  Runtime fuzzing is organized by CPython version, and abnormal outcomes are reported as robustness findings requiring triage rather than automatically as vulnerabilities.

\noindent
{\bf Datasets.}
The package dataset is drawn from PyPI.  For each artifact, we record distribution type, file inventory, source availability, bytecode version, hashes, and dynamic loading indicators.  RQ3 uses CPython unittest-derived bytecode seeds prepared per interpreter version rather than package bytecode, because the robustness question concerns adversarial bytecode processing rather than the benign artifacts observed in RQ1.

\noindent
{\bf Baselines and Tools.}
For RQ2 and RQ4, we use selected bytecode-analysis and decompilation tools, including standard-library loading/disassembly and available decompilers.  Tool failures are recorded with explicit reasons such as missing interpreter, timeout, unsupported bytecode, nonzero exit, or no usable source output.

\noindent
\textbf{Performance Metrics.}
We report package-level prevalence, bytecode transparency, practical analyzability levels, per-tool success and failure reasons, fuzzing outcome distributions, unique abnormal outcomes when available, and tool-bounded source-reproduction classifications.

\noindent
\textbf{Evaluation Sufficiency.}
The evaluation is structured to support the paper's motivation end to end.  RQ1 establishes whether bytecode artifacts appear and whether they are transparent to source-level inspection.  RQ2 tests whether selected tools can practically analyze the observed bytecode in a version-aware setup.  RQ3 measures the runtime boundary exposed by malformed or adversarial bytecode.  RQ4 determines whether selected source-recovery workflows can reproduce bytecode-level findings.  Together, these questions support the central claim only if prevalence, analyzability limits, and robustness findings are interpreted with their stated scope.


\subsection{RQ1: Bytecode Prevalence}\label{sec:effmeasure}

RQ1 measures how often Python bytecode artifacts appear in real distributions.  We count packages containing \texttt{.pyc} files, \texttt{\_\_pycache\_\_} directories, source-less bytecode modules, bundled bytecode artifacts, dynamically loaded code objects, and PyInstaller-style payloads.  Results should be reported both per package and per artifact, because a single package can contain many bytecode objects.

\begin{table}[t]
  \centering
  \caption{Bytecode prevalence across datasets. Replace placeholders with measured values.}
  \label{tab:prevalence}
  \begin{tabular}{lrrrr}
    \toprule
    Dataset & Packages & With \texttt{.pyc} & Source-less & Dynamic load \\
    \midrule
    Popular PyPI & TBD & TBD & TBD & TBD \\
    Random PyPI & TBD & TBD & TBD & TBD \\
    Suspicious/Malicious & TBD & TBD & TBD & TBD \\
    Bundles & TBD & TBD & TBD & TBD \\
    \bottomrule
  \end{tabular}
\end{table}

The main reporting claim for RQ1 should identify whether bytecode artifacts are rare leftovers or recurring distribution artifacts.  We also report common locations, including cache directories, package subtrees, loader resources, archive members, and bundled module stores.


\subsection{RQ2: Practical Analyzability}\label{sec:RQ2}

RQ2 asks whether observed bytecode artifacts can be analyzed by selected tools in a version-aware environment.  For each \texttt{.pyc} file found in RQ1, we run standard-library loading and disassembly, then selected decompilers when available.  Each tool result records both the status and the concrete failure reason, so failures caused by missing interpreters, unsupported bytecode versions, timeouts, and decompiler errors are not conflated.

\begin{table}[t]
  \centering
  \caption{Practical analyzability levels for observed bytecode artifacts. Replace placeholders with measured values.}
  \label{tab:tool-analyzability}
  \begin{tabular}{lrrrrr}
    \toprule
    Version & \texttt{.pyc} files & Loadable & Disassemblable & Partial decompile & Full decompile \\
    \midrule
    CPython 3.8 & TBD & TBD & TBD & TBD & TBD \\
    CPython 3.9 & TBD & TBD & TBD & TBD & TBD \\
    CPython 3.10 & TBD & TBD & TBD & TBD & TBD \\
    CPython 3.12 & TBD & TBD & TBD & TBD & TBD \\
    \bottomrule
  \end{tabular}
\end{table}

The key interpretation is practical rather than theoretical.  If bytecode becomes analyzable after the matching CPython interpreter is installed, the earlier failure is an environment limitation.  If selected tools still fail in a matched environment, the result indicates a current tooling limitation.  RQ2 therefore reports what the selected analysis stack can inspect today, not a proof that the bytecode is or is not recoverable by all possible tools.

\subsection{RQ3: Tooling Gap}\label{sec:RQ3}

RQ3 evaluates how well current tools recover or analyze Python bytecode artifacts.  For each decompiler, we measure decompilation success, syntactic validity of recovered source, recompilation success, bytecode equivalence after recompilation, and behavior preservation when safe execution is possible.  We also test whether source-level scanners observe behavior that exists only in bytecode artifacts.

\begin{table}[t]
  \centering
  \caption{Recovery workflow outcomes. Replace placeholders with measured values.}
  \label{tab:recovery}
  \begin{tabular}{lrrrr}
    \toprule
    Tool & Decompiled & Recompiled & Equivalent & Main failure \\
    \midrule
    pycdc & TBD & TBD & TBD & TBD \\
    Other & TBD & TBD & TBD & TBD \\
    \bottomrule
  \end{tabular}
\end{table}

Failure causes include unsupported Python versions, unsupported opcodes, unusual control flow, exception-table mismatches, malformed code-object metadata, and dynamic loading patterns that hide the bytecode from ordinary file-based analysis.  The main interpretation is conservative: a decompiler failure is a tooling limitation, not proof that the bytecode has no source or no behavior.



\subsection{RQ4: Source Reproduction}\label{sec:eval:runtime}

RQ4 reports whether selected source-recovery workflows can reproduce bytecode-level findings from RQ3.  We classify each finding by the first blocking point in the workflow: no selected decompiler available, decompilation failure, uncompilable source, behavior mismatch, or same-behavior reproduction.

\begin{table}[t]
  \centering
  \caption{Tool-bounded source reproduction for RQ3 findings. Replace placeholders with measured counts.}
  \label{tab:source-reproduction}
  \begin{tabular}{lrrrrr}
    \toprule
    Version & Findings & Decompiled & Compiled & Reproduced & Not reproduced \\
    \midrule
    3.10 & TBD & TBD & TBD & TBD & TBD \\
    3.11 & TBD & TBD & TBD & TBD & TBD \\
    3.12 & TBD & TBD & TBD & TBD & TBD \\
    \bottomrule
  \end{tabular}
\end{table}

For each abnormal outcome, we attempt decompilation, recompilation, and re-execution under the same runtime using the selected tools.  A finding is source-reproduced only if the recovered source produces the same behavior class.  Otherwise, it is classified as not reproduced by selected tools, with the concrete reason recorded as decompilation failed, recompilation failed, behavior changed, tool unavailable, or inconclusive.  This phrasing is deliberate: it avoids claiming that no source reproducer exists when the study only shows that the selected tools did not produce one.

\section{Discussion} \label{sec:diss}

\subsection{Implications for Package Security}

Source-only package analysis is incomplete when distributions contain executable bytecode without matching source.  Package scanners should inventory \texttt{.pyc} files, inspect marshalled code objects, flag bytecode/source mismatches, and preserve decompilation failures as evidence for manual review.  A failed source recovery step should not be treated as a clean result.

\subsection{Implications for Runtime Hardening}

CPython need not support arbitrary malformed bytecode as trusted input, but it should fail predictably at documented validation boundaries.  Assertions, aborts, sanitizer findings, and segmentation faults indicate hardening opportunities even when the triggering artifact is malformed.  The source-reproduction analysis helps separate findings reproduced by selected source workflows from findings that, in this study, require direct bytecode construction.

\subsection{Implications for Malware Analysis}

Compiled Python payloads complicate triage because the analyst may see a small source-level loader while the payload resides in a code object.  Bytecode-aware signatures, opcode-level behavior extraction, and explicit handling for version-specific bytecode can improve malware analysis workflows.  Decompilation remains useful, but it should be one component of a broader bytecode analysis pipeline.

\subsection{Threats to Validity}

Dataset bias is a primary threat: PyPI packages may not represent all Python deployments, including containers, application bundles, or private release artifacts.  Tool bias is also possible because analyzability and source-reproduction results depend on selected tools and versions.  Runtime fuzzing may overproduce malformed inputs that ordinary source would not generate, and finding deduplication can over- or under-count unique failures.  We mitigate these threats by reporting failure categories explicitly, using version-matched interpreters, and separating bytecode robustness from tool-bounded source reproduction.

\subsection{Limitations}

Several limitations remain.  First, {\tech} cannot prove that all dynamically loaded bytecode has been recovered; loaders may decrypt, decompress, download, or synthesize code objects in ways that are not visible statically.  Second, the PyPI dataset does not cover every Python deployment format.  Third, bytecode transparency based on nearby source is conservative and does not by itself establish semantic equivalence.  Fourth, decompiler results reflect the selected tools and versions; future tools may recover artifacts that are currently classified as tool failures.  Fifth, robustness findings from malformed bytecode are not automatically security vulnerabilities.  They require manual triage, source-reproduction analysis, and responsible disclosure when appropriate.

\section{Related Work} \label{sec:related}

\begin{itemize}[leftmargin=1.25em]
  \item Organize prior work into meaningful categories.
  \item For each category, summarize trends and representative work.
  \item Clearly state how your work differs from existing research.
\end{itemize}

\section{Conclusion}\label{sec:conclude}

\begin{itemize}[leftmargin=1.25em]
  \item Reiterate the problem and your main solution.
  \item Summarize key empirical results.
  \item State high-level impact or future research directions.
\end{itemize}



%-------------------------------------------------------------------------------
%\section*{Acknowledgments}
%-------------------------------------------------------------------------------

%The USENIX latex style is old and very tired, which is why
%there's no \textbackslash{}acks command for you to use when
%acknowledging. Sorry.

%\textbf{Do not include any acknowledgements in your submission which may deanonymize you (e.g., because of specific affiliations or grants you acknowledge)}

%-------------------------------------------------------------------------------
% optional clearing of the page
%\cleardoublepage
\appendix


\section{Additional Methodology Details}\label{app:method}

The artifact release should include dataset manifests, package hashes, extraction scripts, interpreter versions, decompiler versions, runtime command lines, timeout settings, and deduplication rules.  For each reported abnormal runtime outcome, the appendix should include the failure class, interpreter build, minimized input hash, top stack frame, source-reachability classification, and disclosure status when applicable.


{
\small
\bibliographystyle{plain}
\bibliography{reference}
}

%%%%%%%%%%%%%%%%%%%%%%%%%%%%%%%%%%%%%%%%%%%%%%%%%%%%%%%%%%%%%%%%%%%%%%%%%%%%%%%%
\end{document}
%%%%%%%%%%%%%%%%%%%%%%%%%%%%%%%%%%%%%%%%%%%%%%%%%%%%%%%%%%%%%%%%%%%%%%%%%%%%%%%%

%%  LocalWords:  endnotes includegraphics fread ptr nobj noindent
%%  LocalWords:  pdflatex acks
